\documentclass[a4paper,12pt]{article}
\usepackage[T2A]{fontenc}
\usepackage[utf8]{inputenc}
\usepackage[russian]{babel}
\usepackage{amsmath,amssymb,mathtools}
\usepackage[sfdefault,type1]{noto}
\usepackage{icomma}
\usepackage[a4paper,left=21mm,right=21mm,top=18mm,bottom=18mm]{geometry}
\usepackage{microtype}
\usepackage{tikz}
\usetikzlibrary{calc,arrows.meta,positioning,shapes.geometric}
\usepackage{pgfplots}
\pgfplotsset{compat=1.18}
\usepackage{booktabs,tabularx}
\usepackage{enumitem}
\usepackage{fancyhdr}
\usepackage{setspace}
\usepackage{xcolor}
\usepackage{hyperref}
\hypersetup{hidelinks,pdftitle={Чек-лист. Сила упругости и закон Гука},pdfauthor={Лёвин Артём Александрович}}

\definecolor{accent}{HTML}{16767A}
\definecolor{darkaccent}{HTML}{225256}
\definecolor{textgray}{HTML}{33393D}
\definecolor{muted}{HTML}{697276}
\color{textgray}
\setstretch{1.17}
\setlength{\parindent}{0pt}
\setlength{\parskip}{2pt}
\setlist[itemize]{leftmargin=7mm,itemsep=1pt,topsep=2pt,parsep=0pt}
\setlist[enumerate]{leftmargin=8mm,itemsep=7pt,topsep=4pt,parsep=0pt}
\newcommand{\key}[1]{\textcolor{darkaccent}{\textbf{#1}}}
\newcommand{\sect}[1]{\vspace{6pt}{\Large\bfseries\color{darkaccent}#1}\par\vspace{3pt}{\color{accent!65}\rule{\linewidth}{0.5pt}}\par\vspace{3pt}}
\newcommand{\subsect}[1]{\vspace{6pt}{\large\bfseries\color{darkaccent}#1}\par\vspace{1pt}}
\newcommand{\tinykey}[1]{\textcolor{accent}{\textbf{#1}}}
\pagestyle{fancy}
\fancyhf{}
\fancyhead[L]{\small\color{darkaccent} Физика \textbullet\ Сила упругости}
\fancyhead[R]{\small\color{darkaccent} Володя Хачатурян}
\fancyfoot[C]{\small\color{muted}\thepage}
\renewcommand{\headrulewidth}{0.35pt}
\renewcommand{\footrulewidth}{0pt}
\setlength{\headheight}{15pt}

\begin{document}
\hypersetup{pageanchor=false}
\begin{titlepage}
\thispagestyle{empty}
\begin{center}
\vspace*{25mm}
{\large\color{darkaccent} ФИЗИКА}\par
\vspace{12mm}
{\fontsize{27}{34}\selectfont\bfseries Чек-лист}\par
\vspace{12mm}
{\fontsize{22}{28}\selectfont\color{accent} Сила упругости\\[3mm]и закон Гука}\par
\vspace{12mm}
{\normalsize Деформация \textbullet\ Равновесие \textbullet\ Графики \textbullet\ Сравнение опытов}\par
\vspace{28mm}
{\Large Володя Хачатурян}\par
\vspace{12mm}
{\large 28 октября 2026 г.}\par
\vfill
{\normalsize Лёвин Артём Александрович}\\[2mm]
{\small эксклюзивно для Володи Хачатуряна}\par
\vspace{20mm}
\end{center}
\begin{tikzpicture}[remember picture,overlay]
\draw[accent,line width=0.95pt]
  ([xshift=36mm,yshift=24mm]current page.south west)
  .. controls ([xshift=65mm,yshift=36mm]current page.south west) and
  ([xshift=93mm,yshift=13mm]current page.south west)
  .. ([xshift=113mm,yshift=25mm]current page.south west)
  .. controls ([xshift=135mm,yshift=37mm]current page.south west) and
  ([xshift=161mm,yshift=12mm]current page.south west)
  .. ([xshift=-36mm,yshift=24mm]current page.south east);
\end{tikzpicture}
\end{titlepage}
\hypersetup{pageanchor=true}

\sect{1. Деформация и сила упругости}
\key{Деформация} --- изменение формы или размеров тела под действием внешних факторов. Возможны растяжение, сжатие, изгиб, сдвиг и кручение.

\begin{itemize}
\item \key{Упругая деформация:} после снятия нагрузки тело восстанавливает исходные форму и размеры.
\item \key{Пластическая деформация:} после снятия нагрузки сохраняется остаточное изменение формы или размеров.
\end{itemize}

\key{Сила упругости} возникает при деформации и стремится вернуть тело к недеформированному состоянию. Закон Гука справедлив в области \key{малых упругих деформаций}; при больших деформациях линейная зависимость может нарушаться.

\subsect{Закон Гука: модуль силы}
\[
\boxed{F_{\text{упр}}=k\,\Delta x},\qquad
\boxed{\Delta x=|l-l_0|}.
\]
Здесь $k$ --- жёсткость тела (для пружины --- жёсткость пружины); $l_0$ --- начальная длина, $l$ --- конечная длина, $\Delta x$ --- \key{модуль изменения длины}, $F_{\text{упр}}$ --- модуль силы упругости.

\begin{center}
\begin{tabular}{@{}lll@{}}
\toprule
Величина & Обозначение & Единица СИ \\
\midrule
Сила & $F_{\text{упр}}$ & ньютон (Н) \\
Жёсткость & $k$ & Н/м \\
Длина и деформация & $l_0,\,l,\,\Delta x$ & метр (м) \\
\bottomrule
\end{tabular}
\end{center}

\subsect{Как определить деформацию}
Пружину длиной $15$~см растянули до $21$~см:
\[
\Delta x=|21-15|=6\ \text{см}=0{,}06\ \text{м}.
\]
Пружину длиной $15$~см сжали до $12$~см:
\[
\Delta x=|12-15|=3\ \text{см}=0{,}03\ \text{м}.
\]

\subsect{Направление силы}
Если $x=l-l_0$ --- изменение длины со знаком, то проекция силы упругости на ось вдоль пружины равна
\[
\boxed{F_{\text{упр},x}=-kx}.
\]
\key{Минус} показывает: сила упругости стремится восстановить исходную длину. При расчёте \emph{модуля} силы используется положительная формула $F_{\text{упр}}=k\Delta x$.

\clearpage
\sect{2. Три основных расчёта}
Из закона Гука можно выразить любую из трёх величин:
\[
\boxed{F_{\text{упр}}=k\Delta x},\qquad
\boxed{k=\frac{F_{\text{упр}}}{\Delta x}},\qquad
\boxed{\Delta x=\frac{F_{\text{упр}}}{k}}.
\]

\subsect{Пример 1. Удлинение по силе и жёсткости}
Сила $F_{\text{упр}}=14$~Н, жёсткость $k=350$~Н/м. Найдём удлинение:
\[
\Delta x=\frac{14}{350}=0{,}04\ \text{м}=4\ \text{см}.
\]
\key{Проверка размерности:}
\[
[\Delta x]=\frac{\text{Н}}{\text{Н/м}}=\text{м}.
\]

\subsect{Пример 2. Одна и та же пружина, два опыта}
При $\Delta x_1=4$~см сила равна $F_1=3$~Н. Какова сила при $\Delta x_2=6$~см?
\[
F_1=k\Delta x_1,\quad F_2=k\Delta x_2
\quad\Longrightarrow\quad
\frac{F_2}{F_1}=\frac{\Delta x_2}{\Delta x_1}.
\]
\[
F_2=3\cdot\frac{6}{4}=\boxed{4{,}5\ \text{Н}}.
\]
\key{Ключевая идея:} если пружина одна и деформации остаются упругими, $k$ не изменяется. При делении двух уравнений общая жёсткость сокращается.

\sect{3. Равновесие подвешенного груза}
На неподвижный груз действуют две силы: вниз --- тяжесть $mg$, вверх --- сила упругости пружины. \key{При равновесии} модули сил равны:
\[
\boxed{k\Delta x=mg},\qquad \boxed{\Delta x=\frac{mg}{k}}.
\]
Массу $m$ подставляем в килограммах; в задачах используем $g=10$~Н/кг, если не указано другое.

\subsect{Пример 3. Добавили одинаковые грузы}
Первоначальная масса $m_1=0{,}1$~кг растягивает пружину на $2{,}5$~см. Добавили ещё три груза по $0{,}1$~кг. Новая масса $m_2=0{,}4$~кг. Тогда
\[
\frac{k\Delta x_2}{k\Delta x_1}=\frac{m_2g}{m_1g}
\quad\Longrightarrow\quad
\Delta x_2=\Delta x_1\frac{m_2}{m_1}
=2{,}5\cdot4=\boxed{10\ \text{см}}.
\]
Этот приём применим, пока пружина подчиняется закону Гука и находится в равновесии.

\clearpage
\sect{4. Как читать график $F(\Delta x)$}
Для одной пружины в области действия закона Гука график модуля силы упругости от модуля удлинения --- прямая через начало координат. \key{Её наклон показывает жёсткость:}
\[
\boxed{k=\frac{F}{\Delta x}}.
\]

\begin{center}
\begin{tikzpicture}
\begin{axis}[
 width=0.83\linewidth,height=62mm,
 xmin=0,xmax=6.5,ymin=0,ymax=16,
 xtick={0,1,2,3,4,5,6},ytick={0,4,8,12,16},
 xlabel={$\Delta x$, см},ylabel={$F_{\text{упр}}$, Н},
 axis lines=left,axis line style={-{Latex[length=2mm]}},
 grid=both,minor tick num=0,major grid style={draw=black!12},
 tick label style={font=\small},label style={font=\small},
 every axis plot/.append style={thick},clip=false]
\addplot[color=accent,domain=0:6,samples=2]{2*x};
\addplot[only marks,mark=*,mark size=2pt,color=darkaccent] coordinates {(4,8)};
\node[anchor=west,font=\small,fill=white,inner sep=2pt] at (axis cs:4.15,8) {$4\ \text{см};\ 8\ \text{Н}$};
\end{axis}
\end{tikzpicture}
\end{center}
Из отмеченной точки: $\Delta x=4$~см $=0{,}04$~м, $F=8$~Н, поэтому
\[
k=\frac{8}{0{,}04}=\boxed{200\ \text{Н/м}}.
\]
Для неподвижного груза массой $2$~кг: $\Delta x=mg/k=20/200=0{,}1$~м $=10$~см.

\subsect{Сравнение двух графиков}
При одном и том же удлинении \key{выше расположена линия более жёсткой пружины}: большему $F$ соответствует больший $k$. Если одинаковы силы, более жёсткая пружина удлиняется меньше.

\sect{5. Повышенный уровень: однородная проволока}
Для \key{прямой однородной проволоки (или стержня)}, растягиваемой вдоль оси в области линейной упругости,
\[
\boxed{k=\frac{ES}{l_0}},\qquad
\boxed{\Delta x=\frac{F l_0}{ES}}.
\]
$E$ --- модуль Юнга (Па $=\text{Н/м}^2$), $S$ --- площадь поперечного сечения (м$^2$), $l_0$ --- исходная длина (м). При неизменных $E,S$: чем длиннее проволока, тем меньше её жёсткость. При неизменных $E,l_0$: чем больше $S$, тем выше жёсткость.

\textcolor{darkaccent}{\textbf{Важно:}} формулу $k=ES/l_0$ нельзя напрямую переносить на обычную витую пружину: для неё жёсткость зависит также от геометрии витков и сдвиговых свойств материала.

\clearpage
\begin{center}
{\LARGE\bfseries\color{darkaccent} Алгоритм: сравнение двух опытов}\par
\vspace{9mm}
{\small Как исключить неизвестную жёсткость и быстро найти новую силу или деформацию}
\end{center}
\vspace{9mm}
\begin{center}
\begin{tikzpicture}[
 >=Latex,node distance=10mm,
 terminal/.style={ellipse,draw=accent,line width=0.9pt,align=center,text width=70mm,minimum height=11mm,inner sep=2mm,font=\small},
 flowstep/.style={rounded corners=2mm,rectangle,draw=accent!85,line width=0.8pt,align=center,text width=116mm,minimum height=19mm,inner sep=3mm,font=\small},
 arrow/.style={-{Latex[length=2.8mm]},draw=darkaccent,line width=0.9pt}]
\node[terminal] (a) {Начало};
\node[flowstep,below=of a] (b) {Определите два состояния: \\ $F_1,\Delta x_1$ и $F_2,\Delta x_2$.\\ Убедитесь, что действует закон Гука.};
\node[flowstep,below=of b] (c) {Запишите \(F_1=k\Delta x_1\) и \(F_2=k\Delta x_2\).\\ Для неподвижного груза можно заменить \(F_i\) на \(m_i g\).};
\node[flowstep,below=of c] (d) {Разделите равенства одно на другое:\\ \(\displaystyle \frac{F_2}{F_1}=\frac{k\Delta x_2}{k\Delta x_1}\).};
\node[flowstep,below=of d] (e) {Сократите общие величины.\\ \(\displaystyle\frac{F_2}{F_1}=\frac{\Delta x_2}{\Delta x_1}\).};
\node[flowstep,below=of e] (f) {Выразите искомое, подставьте данные\\ и проверьте единицы измерения.};
\node[terminal,below=of f] (g) {Ответ и проверка смысла};
\foreach \u/\v in {a/b,b/c,c/d,d/e,e/f,f/g} {\draw[arrow] (\u.south)--(\v.north);}
\end{tikzpicture}
\end{center}
\vspace{5mm}
\textbf{Когда метод меняется:} если сравниваются \emph{разные} пружины, коэффициенты $k_1$ и $k_2$ сокращать нельзя. При одинаковых удлинениях используйте $F_2/F_1=k_2/k_1$.

\clearpage
\sect{6. Типичные ошибки и быстрая проверка}
\begin{itemize}
\item \key{Сантиметры вместо метров:} $4$~см $=0{,}04$~м; $1$~мм $=0{,}001$~м.
\item \key{Путаница между длиной и деформацией:} в закон Гука входит изменение длины $|l-l_0|$, а не конечная длина $l$.
\item \key{Минус в законе Гука:} формула с минусом относится к проекции силы; модуль силы всегда неотрицателен.
\item \key{Сокращение разных жёсткостей:} $k$ сокращается только для одной и той же неизменной пружины.
\item \key{Забытое равновесие:} равенство $k\Delta x=mg$ используется для покоящегося груза, когда эти силы уравновешены.
\item \key{Путаница в пропорциях:} большему растяжению одной пружины соответствует большая сила; при неизменной силе большему $k$ соответствует меньшее растяжение.
\end{itemize}

\subsect{Короткий маршрут решения}
\key{Дано $F,k$?} Выразите $\Delta x=F/k$. \quad
\key{Дано $m,k$ и равновесие?} Используйте $\Delta x=mg/k$.

\key{Даны два опыта?} Запишите два уравнения и разделите. \quad
\key{Дан график?} Сначала найдите $k=F/\Delta x$ по точке графика.

\clearpage
\sect{7. Тренировочный блок. Путь А}
\textcolor{muted}{\small 10 заданий по нарастающей сложности. При вычислениях используйте $g=10$~Н/кг. Если не оговорено иное, деформации упругие и пружины подчиняются закону Гука.}

\begin{enumerate}[label=\textbf{\arabic*.}]
\item Пружина имела длину $16$~см. После растяжения её длина стала $22$~см. Найдите модуль удлинения в сантиметрах и метрах.

\item Пружину длиной $25$~см сжали до $21$~см. Найдите модуль деформации. Укажите направление силы упругости относительно направления сжатия.

\item Жёсткость пружины равна $400$~Н/м, удлинение --- $1{,}5$~см. Найдите модуль силы упругости.

\item Под действием силы $18$~Н пружина удлинилась. Жёсткость пружины $600$~Н/м. Найдите удлинение в сантиметрах.

\item При удлинении одной пружины на $2$~см сила упругости составляет $5$~Н. Какой станет сила при удлинении на $7$~см?

\item Две пружины растянули на одинаковую величину. Первая имеет жёсткость $250$~Н/м, и сила её упругости составляет $10$~Н. Жёсткость второй $400$~Н/м. Найдите силу упругости второй пружины.

\item Неподвижный груз массой $0{,}24$~кг подвешен к вертикальной пружине жёсткостью $120$~Н/м. Найдите удлинение пружины в сантиметрах.

\item Груз массой $0{,}2$~кг растягивает вертикальную пружину на $5$~см. К нему добавили ещё два груза по $0{,}2$~кг каждый. Каково новое удлинение при равновесии?
\end{enumerate}

\clearpage
\sect{Тренировочный блок. Продолжение}
\begin{enumerate}[label=\textbf{\arabic*.},start=9]
\item На графике дана зависимость модуля силы упругости от удлинения одной пружины. Найдите жёсткость пружины и удлинение при подвешивании неподвижного груза массой $1{,}2$~кг.

\begin{center}
\begin{tikzpicture}
\begin{axis}[
 width=0.77\linewidth,height=64mm,
 xmin=0,xmax=7.2,ymin=0,ymax=15,
 xtick={0,1,2,3,4,5,6,7},ytick={0,3,6,9,12,15},
 xlabel={$\Delta x$, см},ylabel={$F_{\text{упр}}$, Н},
 axis lines=left,axis line style={-{Latex[length=2mm]}},
 grid=both,major grid style={draw=black!15},
 tick label style={font=\small},label style={font=\small},clip=false]
\addplot[color=accent,thick,domain=0:7,samples=2]{2*x};
\addplot[only marks,mark=*,mark size=2.1pt,color=darkaccent] coordinates {(4,8)};
\draw[densely dashed,draw=muted] (axis cs:4,0)--(axis cs:4,8)--(axis cs:0,8);
\node[anchor=south east,font=\small,fill=white,inner sep=2pt] at (axis cs:4,8) {$(4;8)$};
\end{axis}
\end{tikzpicture}
\end{center}

\item Прямую однородную проволоку длиной $1{,}2$~м и площадью поперечного сечения $1{,}5$~мм$^2$ растягивают силой $150$~Н. Модуль Юнга материала $E=2\cdot10^{11}$~Па. Найдите абсолютное удлинение в миллиметрах. Затем определите, каким оно станет, если \emph{при той же силе} площадь поперечного сечения увеличить вдвое.
\end{enumerate}

\vspace{6mm}
\textcolor{darkaccent}{\rule{\linewidth}{0.5pt}}
\vspace{3mm}
\textcolor{muted}{\small Для задач 5--8 полезен метод двух уравнений. Для задачи 9 сначала найдите наклон графика с переводом сантиметров в метры. Для задачи 10 используйте модель прямой проволоки.}

\clearpage
\sect{8. Ответы к тренировочному блоку}
\vspace{6mm}
\renewcommand{\arraystretch}{1.65}
\begin{center}
\begin{tabularx}{\linewidth}{@{}cX@{}}
\toprule
№ & Ответ \\
\midrule
1 & $6$~см $=0{,}06$~м. \\
2 & $4$~см $=0{,}04$~м; сила упругости направлена против сжатия. \\
3 & $6$~Н. \\
4 & $3$~см. \\
5 & $17{,}5$~Н. \\
6 & $16$~Н. \\
7 & $2$~см. \\
8 & $15$~см. \\
9 & $k=200$~Н/м; $\Delta x=6$~см. \\
10 & $0{,}6$~мм; после удвоения площади --- $0{,}3$~мм. \\
\bottomrule
\end{tabularx}
\end{center}
\vfill
\begin{center}
{\small\color{muted}Лёвин Артём Александрович --- эксклюзивно для Володи Хачатуряна}\par
\end{center}
\end{document}
